

%%%%%%%%%%%%%%%%%%%%%%%%%%%% Document Setup %%%%%%%%%%%%%%%%%%%%%%%%%%%%

% Don't like 10pt? Try 11pt or 12pt
\documentclass[10pt]{article}

\usepackage[utf8]{inputenc}



% This is a helpful package that puts math inside length specifications
\usepackage{calc}


% bibsection environment removed — now using biblatex for bibliography

% Layout: Puts the section titles on left side of page
\reversemarginpar

%
%         PAPER SIZE, PAGE NUMBER, AND DOCUMENT LAYOUT NOTES:
%
% The next \usepackage line changes the layout for CV style section
% headings as marginal notes. It also sets up the paper size as either
% letter or A4. By default, letter was used. If A4 paper is desired,
% comment out the letterpaper lines and uncomment the a4paper lines.
%
% As you can see, the margin widths and section title widths can be
% easily adjusted.
%
% ALSO: Notice that the includefoot option can be commented OUT in order
% to put the PAGE NUMBER *IN* the bottom margin. This will make the
% effective text area larger.
%
% IF YOU WISH TO REMOVE THE ``of LASTPAGE'' next to each page number,
% see the note about the +LP and -LP lines below. Comment out the +LP
% and uncomment the -LP.
%
% IF YOU WISH TO REMOVE PAGE NUMBERS, be sure that the includefoot line
% is uncommented and ALSO uncomment the \pagestyle{empty} a few lines
% below.
%

%% Use these lines for letter-sized paper
%\usepackage[paper=letterpaper,
            %includefoot, % Uncomment to put page number above margin
%            marginparwidth=1.2in,     % Length of section titles
%            marginparsep=.05in,       % Space between titles and text
%            margin=1in,               % 1 inch margins
%            includemp]{geometry}

%% Use these lines for A4-sized paper
\usepackage[paper=a4paper,
            %includefoot, % Uncomment to put page number above margin
            marginparwidth=30.5mm,    % Length of section titles
            marginparsep=1.5mm,       % Space between titles and text
            margin=25mm,              % 25mm margins
            includemp]{geometry}

%% More layout: Get rid of indenting throughout entire document
\setlength{\parindent}{0in}

%% This gives us fun enumeration environments. compactitem will be nice.
\usepackage{paralist}

%% Reference the last page in the page number
%
% NOTE: comment the +LP line and uncomment the -LP line to have page
%       numbers without the ``of ##'' last page reference)
%
% NOTE: uncomment the \pagestyle{empty} line to get rid of all page
%       numbers (make sure includefoot is commented out above)
%
\usepackage{fancyhdr,lastpage}
\pagestyle{fancy}
%\pagestyle{empty}      % Uncomment this to get rid of page numbers
\fancyhf{}\renewcommand{\headrulewidth}{0pt}
\fancyfootoffset{\marginparsep+\marginparwidth}
\newlength{\footpageshift}
\setlength{\footpageshift}
          {0.5\textwidth+0.5\marginparsep+0.5\marginparwidth-2in}
\lfoot{\hspace{\footpageshift}%
       \parbox{4in}{\, \hfill %
                    \arabic{page} of \protect\pageref*{LastPage} % +LP
%                    \arabic{page}                               % -LP
                    \hfill \,}}

% APA bibliography style via biblatex
\usepackage[style=apa,sorting=ydnt,backend=biber,uniquename=false,uniquelist=false]{biblatex}
\addbibresource{My papers.bib}

% Remove hanging indent from bibliography entries
\setlength{\bibhang}{0pt}

% Add empty line between bibliography entries
\setlength{\bibitemsep}{\baselineskip}

% Do not print annotation field
\DeclareSourcemap{
  \maps[datatype=bibtex]{
    \map{
      \step[fieldset=annotation,null]
    }
  }
}

% Define categories for CV sections
\DeclareBibliographyCategory{journalpub}
\DeclareBibliographyCategory{bookchapter}
\DeclareBibliographyCategory{nonrefereed}

\addtocategory{journalpub}{%
  rennerSmallholderRSPOCertification2024,%
  pavanelloAirconditioningAdaptationCooling2021,%
  steckelDistributionalImpactsCarbon2021a,%
  rennerEffectsEnergyPrice2019b,%
  rennerHouseholdWelfareCO22018,%
  rennerPovertyDistributionalEffects2018b,%
  jakobFeasibleMitigationActions2014b%
}

\addtocategory{bookchapter}{%
  neffSocialEconomicSituation2019%
}

\addtocategory{nonrefereed}{%
  steckelDistributionalImpactsCarbon2021b,%
  layNotParisTrack2016,%
  rennerBrasilienUndMexiko2012%
}

% Finally, give us PDF bookmarks
\usepackage{color,hyperref}
\definecolor{darkblue}{rgb}{0.0,0.0,0.3}
\hypersetup{colorlinks,breaklinks,
            linkcolor=darkblue,urlcolor=darkblue,
            anchorcolor=darkblue,citecolor=darkblue}

%%%%%%%%%%%%%%%%%%%%%%%% End Document Setup %%%%%%%%%%%%%%%%%%%%%%%%%%%%


%%%%%%%%%%%%%%%%%%%%%%%%%%% Helper Commands %%%%%%%%%%%%%%%%%%%%%%%%%%%%

% The title (name) with a horizontal rule under it
%
% Usage: \makeheading{name}
%
% Place at top of document. It should be the first thing.
\newcommand{\makeheading}[1]%
        {\hspace*{-\marginparsep minus \marginparwidth}%
         \begin{minipage}[t]{\textwidth+\marginparwidth+\marginparsep}%
                {\large \bfseries #1}\\[-0.15\baselineskip]%
                 \rule{\columnwidth}{1pt}%
         \end{minipage}}

% The section headings
%
% Usage: \section{section name}
%
% Follow this section IMMEDIATELY with the first line of the section
% text. Do not put whitespace in between. That is, do this:
%
%       \section{My Information}
%       Here is my information.
%
% and NOT this:
%
%       \section{My Information}
%
%       Here is my information.
%
% Otherwise the top of the section header will not line up with the top
% of the section. Of course, using a single comment character (%) on
% empty lines allows for the function of the first example with the
% readability of the second example.
\renewcommand{\section}[2]%
        {\pagebreak[2]\vspace{1.3\baselineskip}%
         \phantomsection\addcontentsline{toc}{section}{#1}%
         \hspace{0in}%
         \marginpar{
         \raggedright \scshape #1}#2}

% An itemize-style list with lots of space between items
\newenvironment{outerlist}[1][\enskip\textbullet]%
        {\begin{itemize}[#1]}{\end{itemize}%
         \vspace{-.6\baselineskip}}

% An environment IDENTICAL to outerlist that has better pre-list spacing
% when used as the first thing in a \section
\newenvironment{lonelist}[1][\enskip\textbullet]%
        {\vspace{-\baselineskip}\begin{list}{#1}{%
        \setlength{\partopsep}{0pt}%
        \setlength{\topsep}{0pt}}}
        {\end{list}\vspace{-.6\baselineskip}}

% An itemize-style list with little space between items
\newenvironment{innerlist}[1][\enskip\textbullet]%
        {\begin{compactitem}[#1]}{\end{compactitem}}

% An environment IDENTICAL to innerlist that has better pre-list spacing
% when used as the first thing in a \section
\newenvironment{loneinnerlist}[1][\enskip\textbullet]%
        {\vspace{-\baselineskip}\begin{compactitem}[#1]}
        {\end{compactitem}\vspace{-.6\baselineskip}}

% To add some paragraph space between lines.
% This also tells LaTeX to preferably break a page on one of these gaps
% if there is a needed pagebreak nearby.
\newcommand{\blankline}{\quad\pagebreak[2]}

% Uses hyperref to link DOI
\newcommand\doilink[1]{\href{http://dx.doi.org/#1}{#1}}
\newcommand\doi[1]{doi:\doilink{#1}}

% For \url{SOME_URL}, links SOME_URL to the url SOME_URL
%\providecommand*\url[1]{\href{#1}{#1}}
% Same as above, but pretty-prints SOME_URL in teletype fixed-width font
\renewcommand*\url[1]{\href{#1}{\texttt{#1}}}

% For \email{ADDRESS}, links ADDRESS to the url mailto:ADDRESS
\providecommand*\email[1]{\href{mailto:#1}{#1}}
% Same as above, but pretty-prints ADDRESS in teletype fixed-width font
%\renewcommand*\email[1]{\href{mailto:#1}{\texttt{#1}}}

%%%%%%%%%%%%%%%%%%%%%%%% End Helper Commands %%%%%%%%%%%%%%%%%%%%%%%%%%%

%%%%%%%%%%%%%%%%%%%%%%%%% Begin CV Document %%%%%%%%%%%%%%%%%%%%%%%%%%%%

\begin{document}
\makeheading{SEBASTIAN RENNER}

\section{Contact information}
%Mercator Research Institute on Global Commons and Climate Change (MCC)\\
%
% NOTE: Mind where the & separators and \\ breaks are in the following
%       table.
%
% ALSO: \rcollength is the width of the right column of the table
%       (adjust it to your liking; default is 1.85in).
%
\newlength{\rcollength}\setlength{\rcollength}{2.85in}%
%
\begin{tabular}[t]{@{}p{\textwidth-\rcollength}p{\rcollength}}
%\href{http://www.http://www.giga-hamburg.de/}%
  {University of Hamburg} & \\
	{Department of Economics} & \\
%\href{http://www..../}{University of Göttingen}
Von-Melle-Park 5     							 &  \\ %\textit{Mobile:} +49-176-81182009 \\
20146 Hamburg, Germany             & \textit{E-mail:} \email{sebastian.renner@uni-hamburg.de}\\
                             &   \\                                                   
\end{tabular}
Website: \url{https://sebastian-renner.github.io} \\
%ORCID:  \url{https://orcid.org/0000-0002-9448-3386} \\
%h-index: 8 

\section{Research interests}
%
Environmental and Development Economics \& Policy, Science-Policy Interface, Sustainability Transformations, Climate Economics, Causal Inference, Replication Studies 

\section{Education}
%
%\href{http://www.uni-goettingen.de/en/2165.html}{\textbf{University of Göttingen}}\\
\textbf{University of Goettingen}\\
\emph{Ph.D. Economics (Summa cum laude), 2017}\\
Poverty, Inequality and the Decarbonization of Economic Development\\
%\emph{Supervisors: Prof. Stephan Klasen, Prof. Jann Lay, Prof. Krisztina Kis-Katos}\\

\textbf{Free University Berlin}\\
%\emph{Diplom Volkswirt (M.A. Economics), 2010}\\
\emph{M.A. / Diplom in Economics, 2010}\\

\textbf{University of California, San Diego}\\
\emph{Visiting Student, 2007/2008}\\


\section{Current position}
\textbf{University of Hamburg}\\
\emph{Research Associate, since 04/2023}\\
Principal Investigator: Multiple Impacts of Energy Efficiency School Rehabilitation in Georgia (MultEER); with KfW and EBRD, funded by BMZ \\

%Investigator: Robustness and Reproducibility in Economics (R2E); RWI and Stockholm School of Economics \\
%

\section{Previous positions}
\textbf{University of Hamburg}\\
\emph{Visiting Professor (Vertretungsprofessor), 10/2024-03/2025}\\
Teaching: Sustainability Economics, Empirical Environmental Economics, Economics of Climate Change \\

\textbf{Mercator Research Institute on Global Commons and Climate Change (MCC)}\\
\emph{Researcher, 2020-2023}\\
Mobilizing Domestic Resources for Sustainable Development in Indonesia, funded by BMZ\\
Carbon Pricing Design, Economic and Distributional Impact Analysis for the Fiscal Policy Unit (BKF), Ministry of Finance of Indonesia\\ 

\textbf{University of Goettingen}/\textbf{GIGA}\\
\emph{Postdoctoral Researcher, 2016-2020}\\
DFG Collaborative Research Center: Ecological and Socioeconomic Functions of Tropical Lowland Rainforest Transformation Systems (Sumatra, Indonesia)\\

\textbf{German Institute for Global and Area Studies (GIGA) }\\
\emph{Research Fellow, 2011-2020}\\
Climate Change Mitigation and Poverty Reduction, funded by Volkswagen Stiftung, with University of Cape Town (South Africa), National Institute of Development Administration (Thailand), Chiang Mai University (Thailand), Monterrey Institute of Technology and Higher Education (Mexico) and Fondazione Eni Enrico Mattei (Italy) \\
Decarbonization in Low- and Middle Income Countries, funded by BMBF, with PIK and University of Goettingen \\

\section{Affiliations}
\textbf{Environment for Development (EfD) }\\
\emph{International Associate, since 05/2020}\\
\emph{Emission Pricing for Development (EPFD) }\\

\textbf{German Institute for Global and Area Studies (GIGA) }\\
\emph{Associate Research Fellow , since 05/2020}\\


%\newpage 

\section{Working Papers}
Coal-fired Power Plants and Industrial Development, R\&R at Journal of Development Economics \\
with Leonard Missbach, Jan Steckel and Sebastian Kraus\\

Energy Subsidy Reform and Manufacturing Firm Performance: Evidence from Indonesia \\
with Krisztina Kis-Katos\\

Energy Prices, Generators, and the Performance of Manufacturing Firms, with Hannes Greve and Krisztina Kis-Katos\\

Data Gaps, Targeting Transfers and the Incidence of Carbon Pricing in Indonesia, with Djoni Hartono, Jan Steckel and Arief A. Yusuf\\


\section{Work in progress}
Shock Propagation in Global Supply Chains; with Bastian Westbrock and Maarten Bosker \\

EFForTS-ABM v2.0: An Agent-Based Model of the Spatio-Temporal Dynamics of Socio-Economic and Ecological Functions in Oil Palm-Dominated Tropical Landscapes, with Sebastian Fiedler, Julia Henzler, Kerstin Wiegand and Jann Lay\\

%Comment on Akresh et al. (2023, The Economic Journal), Long-Term and Intergenerational Effects of Education: Evidence from School Construction in Indonesia; Institute for Replication (I4R) \\

\section{Journal Publications}
\nocite{rennerSmallholderRSPOCertification2024,pavanelloAirconditioningAdaptationCooling2021,steckelDistributionalImpactsCarbon2021a,rennerEffectsEnergyPrice2019b,rennerHouseholdWelfareCO22018,rennerPovertyDistributionalEffects2018b,jakobFeasibleMitigationActions2014b}
\printbibliography[category=journalpub,heading=none]

\section{Book chapters}
\nocite{neffSocialEconomicSituation2019}
\printbibliography[category=bookchapter,heading=none]

\section{Non-Refereed Publications \& Policy Briefs}
\nocite{steckelDistributionalImpactsCarbon2021b,layNotParisTrack2016,rennerBrasilienUndMexiko2012}
\printbibliography[category=nonrefereed,heading=none]


\section{Professional Experience}

\textbf{United Nations University World Institute for Development Economics Research (UNU-WIDER)}\\
\emph{Investigator, 2026}\\
Taxes, subsidies, and the environment project, part of UNU-WIDER Domestic Revenue Mobilization (DRM) programme\\

\textbf{RWI / Institute for Replication}\\
\emph{Investigator/Replicator, 2024/2025/2026}\\
Robustness and Reproducibility in Economics (R2E)\\

\textbf{KfW Development Bank}\\
\emph{Short Term Consultant, 2020}\\
Ex-post Evaluation of Energy Efficiency Infrastructure Project in Bangladesh\\

\textbf{Agence Française de Développement (AFD)}\\
\emph{Short Term Consultant, 2019}\\
Managing Distributional Impacts of Carbon Taxes in Colombia\\

\textbf{Euro-Mediterranean Center on Climate Change (CMCC)} and \textbf{Ca'Foscari University of Venice}\\
\emph{Consultant, 2018-2020}\\
ENERGYA - Energy Use For Adaptation\\

\textbf{World Bank}\\
\emph{Short Term Consultant, April - June 2013}\\
Development Economics Prospects Group\\

\emph{Short Term Consultant, June 2012}\\
Economic Policy Group, Latin America and the Caribbean Region\\

\emph{Short Term Consultant, February - April 2012}\\
Development Economics Prospects Group\\

\textbf{BC3 Basque Centre for Climate Change}\\
\emph{Visiting Researcher, September - November 2012}\\


\section{Recent Teaching}
Sustainability Economics\\
\emph{University of Hamburg -- Winter 2024/25}\\

Empirical Environmental Economics\\
\emph{University of Hamburg -- Winter 2024/25}\\

Economics of Climate Change\\
\emph{University of Hamburg -- Winter 2024/25}\\


\section{Research Grants}
%DFG Research Grant (submitted, under review): Climate Policy Shocks and Production Networks in the EU-ETS (Spain) \\

%Grant proposal (finalized): Energy Policy Shocks and Production Networks in Middle-Income Countries (Turkey, Ecuador) \\

DEval Rigorous Impact Evaluation / Federal Ministry for Economic Cooperation and Development (BMZ), Multiple impacts of energy efficiency school rehabilitation in Georgia (MultEER) \\

Federal Ministry of Education and Research Germany (BMBF), De-Carbonizing Economic Development in Sub-Saharan Africa (DeCaDe), 2018, with Jan Steckel, Michael Jakob, Jann Lay, Joerg Peters\\

\section{Professional activities}
Referee: Journal of Development Economics, Journal of Environmental Economics and Management, Nature Climate Change, Nature Sustainability, Review of Income and Wealth, Energy Economics, Nature Aging, Environment and Development Economics, Energy Policy, Climate Policy, Energy Research \& Social Science, Agricultural and Food Economics, One Earth, Energy Efficiency, Technological Forecasting \& Social Change, Inter-American Development Bank Working Paper Series, Utilities Policy \\
%Supervisor for student competition YES! � Young Economic Summit 2018 
%Associate Editor Indonesian Journal of Energy\\

%Managing "`GIGA Seminar in Socio-Economics"'\\
Thesis supervision: Master theses in MSc Development Economics and MSc International Economics (12x), University of Göttingen; Bachelor and Master theses at University of Hamburg, integrated into ongoing research projects\\

 
%\section{Languages}
%
%German (native), English (fluent), Spanish (elementary), French (elementary), Bahasa Indonesia (basic)\\

%Computing: Stata, R, Matlab, GAMS, CsPro, {\LaTeX}

%\newpage

%\section{World Bank reference available to contact}
%Maurizio Bussolo, PhD\\
%Lead Economist, Chief Economist Office for South Asia\\
%The World Bank\\
%1818 H Street\\
%NW Washington, DC 20433 USA\\
%E-mail: mbussolo@worldbank.org\\

\end{document}


%%%%%%%%%%%%%%%%%%%%%%%%%% End CV Document %%%%%%%%%%%%%%%%%%%%%%%%%%%%%
